\section{Conclusion}
\label{sec:conclusion}

We introduced boundary regret as the deployment estimand for human--AI complementarity, shifting the audit from ``does the human have residual predictive information?'' to ``does that information change the reward-relevant action?''
The formal characterization---VoI equals expected boundary regret, with a Lean-verified geometric bridge---provides a principled basis for this shift.
Empirically, on CheXpert across five readers and four conditions (20 pairs), residual expertise and decision value diverge in a reward-dependent manner precisely as the theory predicts: a mean of 95.9\% of top-expertise-decile cases have zero decision value under a false-positive-penalty reward.
In review allocation on both synthetic and real CheXpert data, boundary regret outperforms standard baselines under asymmetric rewards while matching them under symmetric accuracy.

\paragraph{Limitations.}
Several limitations frame the contribution.
\emph{Retrospective allocation:} Experiment~3 validates the allocation principle on both synthetic and real CheXpert data, but both evaluations are retrospective.
Prospective evaluation with operational review protocols is required before deployment conclusions can be drawn; the L2D baseline is also trained on accuracy rather than reward-specific utility, which disadvantages it under asymmetric rewards.
\emph{Dataset scope:} Experiment~2 covers five readers and four conditions (20 pairs) on a single dataset (CheXpert).
Generalization to additional datasets, domains, and human-signal types beyond radiology is needed for a broad empirical claim.
\emph{Potential vs.\ realized value:} boundary regret measures the decision value available in the joint distribution $(x, h, y)$, not whether a deployed human--AI workflow can extract it.
Behavioral frictions---anchoring, overreliance, noncompliance---can prevent realization \citep{Vaccaro2024,BucinaEtAl2021}, and bridging from potential to realized value requires prospective studies beyond this paper.
\emph{Estimation quality:} the ex ante estimator $\widehat\BR_{\mathrm{pre}}$ depends on a model of $P(h \mid x)$ whose quality directly affects allocation performance; poor conditional models will degrade the score, and the sample complexity of this estimator is not formally characterized.

\paragraph{Future work.}
Three directions follow naturally.
(1)~Prospective allocation experiments on clinical data, measuring whether boundary-regret-guided review improves realized decision utility under operational conditions and realistic human response models.
(2)~Integration with learning-to-defer methods: boundary regret can serve as an interpretable diagnostic for understanding what learned deferral policies exploit---reward-local boundary crossings or other dataset correlations.
When L2D outperforms boundary regret, the diagnostic can reveal whether the learned policy is capturing true decision-boundary information or overfitting to non-geometric signals.
(3)~Extension to the multi-channel interaction identity (ASIL, Appendix~\ref{app:asil}), which decomposes the complement/substitute relationship between model and human signals into facet-specific contributions, enabling a richer theory of when acquiring a second expert opinion adds value beyond the first.

\paragraph{Code and data availability.}
The Lean~4 formalization and experiment code will be released upon publication. CheXpert data are publicly available \citep{Irvin2019}.
